\section{Introduction}

This paper describes how we collect real-world demonstrations on a dual
five-degree-of-freedom manipulation rig so that the data can train a
policy which runs autonomously on the same platform. It first presents the
collection platform, then the recording system, then the contract that
links a teleoperated demonstration to autonomous control, then the
measurements needed to reproduce the setup, and finally an honest account
of the difficulties met. Each part forward-links to the next.

Behaviour cloning learns a map from an observation at one moment to the
intended command at that same moment~\citep{ross2011dagger}. A
demonstration is therefore only as good as two properties: the observation and the intent must belong to the
same instant, and the recorded intent must be the command the robot
actually executed. Break either property and the policy learns a jittery
or non-causal map that fails when it drives the arm on its own.

Our rig makes both properties hard to hold. Several cameras stream over
shared buses and arrive at slightly different times, while the arm's joint
signals arrive several times faster. If the recorder simply takes each
stream's most recent value when it writes a frame, the streams in that
frame are silently skewed against one another. Separately, the operator's
raw input is not what the arm followed: a rate limiter, the joint limits, a
deliberately capped gripper opening, and the reachable-workspace projection
all sit between the operator and the motors.

The first idea addresses timing. Every stream is timestamped at the instant
it is read from its device, on a monotonic clock. For each recorded frame
the collector fixes a single reference time and samples every stream at
that time, and it records the residual gap between the reference time and
each stream's nearest real sample. Temporal alignment thus becomes a
measured quantity, shown live during collection, rather than an assumption.

The second idea addresses intent. We record the projected and constrained
target --- the command the robot actually followed --- and we record it in
both joint and end-effector form. We then require that, at inference, a
policy's output passes back through the same command path that produced the
labels. A joint-space policy replays through the same limiter and gripper
cap; an end-effector policy replays through the same inverse kinematics and
envelope projection.

Our contributions are: a timestamp-aligned multi-view recorder with a
logged per-stream drift budget (Section~\ref{sec:collection}); a
teleoperation-to-autonomy contract that records the executed, constrained
target for both a joint-space and an end-effector policy
(Section~\ref{sec:contract}); a measured, reproducible platform whose
defining constants travel with the dataset (Sections~\ref{sec:platform}
and~\ref{sec:replicability}); and a candid difficulties record
(Section~\ref{sec:difficulties}).
